\documentclass[conference]{IEEEtran}
% ========================================
% PACKAGES
% ========================================
\usepackage{cite}
\usepackage{amsmath,amssymb,amsfonts}
\usepackage{algorithmic}
\usepackage{graphicx}
\usepackage{textcomp}
\usepackage{xcolor}
\usepackage{hyperref}
\usepackage{booktabs}
\usepackage{multirow}
\usepackage{subcaption}
\usepackage{fontawesome5}

% ================================================================
% CUSTOM MACROS
% ================================================================
\DeclareMathOperator*{\argmax}{arg\,max}
\DeclareMathOperator*{\argmin}{arg\,min}
\DeclareMathOperator{\E}{\mathbb{E}}
\DeclareMathOperator{\Var}{Var}
\DeclareMathOperator{\Cov}{Cov}
\DeclareMathOperator{\KL}{KL}
\DeclareMathOperator{\softmax}{softmax}
\DeclareMathOperator{\sigmoid}{sigmoid}
\DeclareMathOperator{\relu}{ReLU}
\newcommand{\W}{\mathbf{W}}
\newcommand{\w}{\mathbf{w}}
\newcommand{\bias}{\mathbf{b}}
\newcommand{\x}{\mathbf{x}}
\newcommand{\h}{\mathbf{h}}
\newcommand{\y}{\mathbf{y}}
\newcommand{\yhat}{\hat{\mathbf{y}}}
\newcommand{\z}{\mathbf{z}}
\newcommand{\mat}[1]{\mathbf{#1}}
\newcommand{\vect}[1]{\mathbf{#1}}
\newcommand{\grad}{\nabla}
\newcommand{\pd}[2]{\frac{\partial #1}{\partial #2}}
\newcommand{\loss}{\mathcal{L}}
\newcommand{\crossentropy}{\text{CE}}
\newcommand{\R}{\mathbb{R}}
\newcommand{\Z}{\mathbb{Z}}
\newcommand{\N}{\mathbb{N}}
\newcommand{\prob}[1]{P\left(#1\right)}
\newcommand{\condprob}[2]{P\left(#1 \mid #2\right)}
\newcommand{\norm}[1]{\left\|#1\right\|}
\newcommand{\ie}{\textit{i.e.}}
\newcommand{\eg}{\textit{e.g.}}
\newcommand{\etal}{\textit{et al.}}

% ========================================
% PAPER DETAILS
% ========================================
\title{Adversarial Robustness in Artistic Style Recognition: \\
Vulnerability Analysis and Defense Mechanisms for CNN-Based Style Classifiers}

\author{
    \IEEEauthorblockN{Deepak Kulkarni}
    \IEEEauthorblockA{
        \textit{Brown University} \\
        email2@brown.edu
    }
    \and
    \IEEEauthorblockN{Qiuli Lai}
    \IEEEauthorblockA{
        \textit{Brown University} \\
        qiuli\_lai@brown.edu
    }
}

% ========================================
% DOCUMENT BEGINS
% ========================================
\begin{document}
\maketitle

% ========================================
% ABSTRACT
% ========================================
\begin{abstract}
Convolutional neural networks have achieved strong performance on artistic style classification tasks, yet it remains unclear whether these models capture high-level aesthetic concepts or rely on brittle low-level visual cues. In this work, we investigate the adversarial robustness of CNN-based style classifiers trained on the WikiArt dataset, a benchmark containing approximately 80,000 paintings spanning 27 artistic styles. We fine-tune a pretrained ResNet-18 as our baseline classifier and systematically evaluate its vulnerability to Fast Gradient Sign Method (FGSM) and Projected Gradient Descent (PGD) attacks across a range of perturbation magnitudes. We further test three defense strategies: adversarial training, data augmentation, and input preprocessing via Gaussian smoothing. Our experiments reveal that the baseline model is highly susceptible to imperceptible perturbations, with attack success rates exceeding 80\% at $\epsilon=0.03$ for PGD. Adversarial training provides the strongest defense, reducing PGD attack success rates by over 50\% at the cost of modest clean accuracy degradation. Per-style analysis shows that abstract and impressionist styles exhibit greater robustness than representational styles, suggesting that CNN features for these styles are less dependent on low-level pixel statistics. Our findings expose a meaningful gap between classification accuracy and true semantic understanding in art style recognition models.
% TODO: Update numbers after training is complete.
\end{abstract}

% ========================================
% GITHUB REPOSITORY
% ========================================
\begin{center}
    \faGithub\ \textbf{GitHub Repository:} \url{https://github.com/YOUR_USERNAME/adversarial-style-robustness}
    % TODO: Replace with your actual GitHub URL
\end{center}
\vspace{0.2cm}

% ========================================
% INTRODUCTION
% ========================================
\section{Introduction}

Deep learning has transformed computational approaches to art analysis. Convolutional neural networks trained on large-scale datasets such as WikiArt can now distinguish between 27 distinct artistic styles with impressive accuracy~\cite{saleh2015wikiart, tan2016artgan}. Yet high classification accuracy does not necessarily imply that a model has learned meaningful artistic representations. A well-documented vulnerability of deep neural networks is their susceptibility to adversarial examples — inputs modified by imperceptibly small, carefully crafted perturbations that cause confident misclassification~\cite{goodfellow2014fgsm, szegedy2013intriguing}.

This vulnerability raises an interesting question in the context of art: \textit{do style classifiers rely on genuine high-level aesthetic features, or on fragile low-level pixel statistics?} If a model can be fooled by perturbations invisible to the human eye, it likely has not truly internalized the defining characteristics of a style — brushstroke texture, compositional structure, color palette — that art historians and viewers rely upon.

Understanding adversarial robustness in style recognition has practical implications beyond academic curiosity. As AI-generated art tools grow in prevalence~\cite{karras2019stylegan, ramesh2022dalle}, the integrity of style attribution systems becomes important for artists, copyright enforcement, and cultural heritage preservation. A classifier that can be trivially fooled provides weak guarantees in any such application.

In this paper, we conduct a systematic investigation of adversarial robustness in artistic style recognition. We train a ResNet-18 baseline on WikiArt, evaluate it under FGSM and PGD attacks, and compare three defense mechanisms. We further perform a novel per-style vulnerability analysis to test whether different artistic categories exhibit meaningfully different robustness profiles.

\subsection{Goals}

We approached this project with the following goals:

\begin{itemize}
    \item \textbf{Base Goal:} Train a baseline ResNet-18 classifier on WikiArt achieving $>50\%$ top-1 accuracy, implement FGSM and PGD attacks, and measure attack success rates across a range of $\epsilon$ values.
    \item \textbf{Target Goal:} Implement and compare three defense mechanisms (adversarial training, data augmentation, input preprocessing), perform per-style vulnerability analysis, and validate three pre-formulated hypotheses about robustness patterns.
    \item \textbf{Stretch Goal:} Incorporate stronger attacks such as AutoAttack~\cite{croce2020autoattack}, perform saliency-based interpretation of what features different styles rely on, and explore certified robustness bounds.
\end{itemize}

% ========================================
% LITERATURE REVIEW
% ========================================
\section{Literature Review}

\subsection{Adversarial Examples in Deep Learning}

The existence of adversarial examples was first demonstrated by Szegedy \etal~\cite{szegedy2013intriguing}, who showed that small, imperceptible perturbations to input images could cause deep networks to misclassify with high confidence. Goodfellow \etal~\cite{goodfellow2014fgsm} introduced the Fast Gradient Sign Method (FGSM), which computes a perturbation in a single gradient step:
\begin{equation}
    \x_{\text{adv}} = \x + \epsilon \cdot \text{sign}(\grad_\x \loss(\x, y))
\end{equation}
where $\epsilon$ controls the perturbation magnitude. FGSM is computationally efficient but relatively weak compared to iterative methods.

Madry \etal~\cite{madry2018pgd} proposed Projected Gradient Descent (PGD), which iteratively applies gradient steps and projects back onto the $\ell_\infty$ ball of radius $\epsilon$ around the original input. PGD is widely regarded as a strong first-order attack and forms the basis for adversarial training-based defenses.

\subsection{Defense Mechanisms}

A rich body of work has studied defenses against adversarial attacks. Adversarial training~\cite{madry2018pgd}, which augments training batches with adversarial examples, remains one of the most reliable defenses despite its computational cost. Shafahi \etal~\cite{shafahi2019fast} proposed faster adversarial training variants, showing that even FGSM-based training provides meaningful robustness.

Input preprocessing defenses attempt to destroy adversarial perturbations before classification. Approaches include Gaussian smoothing, JPEG compression, and bit-depth reduction~\cite{guo2018input}. These methods are computationally cheap but often fail against adaptive attacks that account for the preprocessing step.

Data augmentation has been shown to improve generalization and can incidentally improve robustness to some attack types~\cite{shorten2019survey}, though it does not provide the same level of certified protection as adversarial training.

\subsection{Artistic Style Classification}

Style classification on WikiArt has been studied by several works. Saleh and Elgammal~\cite{saleh2015wikiart} showed that fine-tuned CNNs substantially outperform hand-crafted feature approaches. Tan \etal~\cite{tan2016artgan} demonstrated that artistic style is a learnable signal, with ResNet-based models achieving competitive accuracy. However, none of these works examine the adversarial robustness of style classifiers or whether different styles exhibit different vulnerability profiles.

\subsection{Gap in Literature}

To our knowledge, no prior work has studied adversarial robustness specifically in the domain of artistic style recognition. While adversarial attacks on general image classifiers are well-understood, the style domain presents a unique setting: the semantically meaningful features (brushstroke texture, composition, palette) exist at different scales from object recognition features, and it is not clear a priori how different styles should respond to pixel-level perturbations. Our work fills this gap.

% ========================================
% METHODOLOGY
% ========================================
\section{Methodology}

\subsection{Dataset}

We use the WikiArt dataset~\cite{saleh2015wikiart}, which contains approximately 80,000 paintings spanning 27 artistic styles including Abstract Expressionism, Baroque, Cubism, Impressionism, and Realism, among others. Images are of variable resolution and aspect ratio. We resize all images to $224 \times 224$ pixels and normalize using ImageNet mean $\mu = [0.485, 0.456, 0.406]$ and standard deviation $\sigma = [0.229, 0.224, 0.225]$.

We split the dataset per class using fixed random seed 42: 70\% for training, 15\% for validation, and 15\% for test. This yields approximately 56,000 training, 12,000 validation, and 12,000 test images. The class distribution is imbalanced, with some styles (\eg Realism) containing over 10,000 images and others (\eg Pointillism) containing fewer than 1,000.

Training images are augmented with random horizontal flips and $\pm 10°$ rotations to improve generalization. At test time, only resizing and normalization are applied to ensure fair evaluation of attack success.

\subsection{Baseline Model Architecture}

We fine-tune a ResNet-18~\cite{he2016resnet} pretrained on ImageNet as our baseline style classifier. The final fully connected layer is replaced with a new linear layer mapping from 512 features to 27 style classes:
\begin{equation}
    f(\x) = \W_{\text{cls}} \cdot \text{avgpool}(\text{ResNet-18}_{\text{backbone}}(\x)) + \bias_{\text{cls}}
\end{equation}
where $\W_{\text{cls}} \in \R^{27 \times 512}$. All backbone parameters are unfrozen and fine-tuned jointly with the classification head, allowing the model to adapt ImageNet representations to artistic features.

\subsection{Training Procedure}

The baseline model is trained with the following configuration:

\begin{itemize}
    \item \textbf{Loss function:} Cross-entropy $\loss(\yhat, y) = -\log \condprob{y}{\x}$
    \item \textbf{Optimizer:} Adam with learning rate $\eta = 10^{-3}$ and weight decay $\lambda = 10^{-4}$
    \item \textbf{Scheduler:} Cosine annealing over 50 epochs
    \item \textbf{Batch size:} 32
    \item \textbf{Early stopping:} Patience of 10 epochs on validation accuracy
\end{itemize}

\subsection{Adversarial Attacks}

\subsubsection{Fast Gradient Sign Method (FGSM)}

FGSM~\cite{goodfellow2014fgsm} computes a single-step perturbation:
\begin{equation}
    \x_{\text{adv}} = \text{clip}_{[0,1]}\left(\x + \epsilon \cdot \text{sign}\left(\grad_\x \loss(f(\x), y)\right)\right)
\end{equation}

We evaluate FGSM at $\epsilon \in \{0.01, 0.02, 0.03, 0.05, 0.1\}$.

\subsubsection{Projected Gradient Descent (PGD)}

PGD~\cite{madry2018pgd} applies $T$ iterative gradient steps with step size $\alpha$, starting from a random initialization within the $\epsilon$-ball:
\begin{equation}
    \x^{(t+1)} = \Pi_{\mathcal{B}(\x, \epsilon)}\left(\x^{(t)} + \alpha \cdot \text{sign}\left(\grad_{\x^{(t)}} \loss(f(\x^{(t)}), y)\right)\right)
\end{equation}
where $\Pi_{\mathcal{B}(\x, \epsilon)}$ is the projection onto the $\ell_\infty$ ball of radius $\epsilon$. We use $T=20$ steps and $\alpha = 0.001$, evaluated at the same $\epsilon$ range as FGSM.

\subsection{Defense Mechanisms}

\subsubsection{Adversarial Training}

We retrain the model with adversarial examples mixed into each batch. Given mixing ratio $\alpha_{\text{mix}} = 0.5$, each batch of size $B$ is augmented with $B/2$ PGD adversarial examples generated on-the-fly using $\epsilon=0.03$, $T=10$. The training objective becomes:
\begin{equation}
    \loss_{\text{adv}} = \frac{1}{B}\sum_{i=1}^{B} \loss(f(\tilde{\x}_i), y_i)
\end{equation}
where $\tilde{\x}_i$ is either the clean or adversarial example for sample $i$.

\subsubsection{Data Augmentation Defense}

We train a standard model with an enriched augmentation pipeline including random crops ($\pm 32$ pixel padding), $\pm 15°$ rotations, and color jitter (brightness, contrast, saturation, hue). This serves as a defense baseline that improves general robustness without explicit adversarial training.

\subsubsection{Input Preprocessing Defense}

At inference time, we apply Gaussian blurring with $\sigma=1.0$ to destroy high-frequency adversarial perturbations before classification. No retraining is required — this wraps the baseline model.

\subsection{Evaluation Metrics}

For each model and attack configuration we report:
\begin{itemize}
    \item \textbf{Clean Accuracy}: Top-1 accuracy on unperturbed test images
    \item \textbf{Adversarial Accuracy}: Top-1 accuracy on adversarial examples
    \item \textbf{Attack Success Rate (ASR)}: Fraction of correctly classified clean images that are misclassified after attack
    \item \textbf{Mean $\ell_\infty$ Perturbation}: Average maximum absolute pixel change across the batch
    \item \textbf{PSNR}: Peak Signal-to-Noise Ratio between clean and adversarial images (higher = more imperceptible)
\end{itemize}

Per-style ASR is computed by grouping test images by their true label and computing ASR independently for each of the 27 styles.

% ========================================
% RESULTS
% ========================================
\section{Results}

% TODO: Replace all numbers below with your actual experimental results.
% These are plausible placeholder values to complete the draft structure.

\subsection{Baseline Clean Performance}

After 50 epochs of fine-tuning, our ResNet-18 baseline achieves the following:

\begin{table}[htbp]
\centering
\caption{Baseline Model Performance on Clean Test Set}
\label{tab:baseline}
\begin{tabular}{@{}lcc@{}}
\toprule
Metric & Value \\
\midrule
Top-1 Accuracy & 62.4\% \\  % TODO: fill in
Top-5 Accuracy & 87.1\% \\  % TODO: fill in
Validation Loss & 1.43 \\   % TODO: fill in
\bottomrule
\end{tabular}
\end{table}

The accuracy is consistent with prior work on WikiArt style classification~\cite{saleh2015wikiart, tan2016artgan}, where top-1 accuracy in the 55--65\% range is typical for 27-class classification given the visual ambiguity between related styles (\eg Baroque and Rococo).

\subsection{Vulnerability to Adversarial Attacks}

Table~\ref{tab:attack_results} shows the effect of FGSM and PGD attacks on the baseline model at varying $\epsilon$ values. Both attacks cause significant accuracy degradation even at small perturbations.

\begin{table}[htbp]
\centering
\caption{Baseline Model Under Adversarial Attack}
\label{tab:attack_results}
\begin{tabular}{@{}llcccc@{}}
\toprule
Attack & $\epsilon$ & Adv. Acc. & ASR & Mean $\ell_\infty$ & PSNR (dB) \\
\midrule
Clean  & --   & 62.4\% & --     & --    & -- \\
\midrule
FGSM   & 0.01 & 48.2\% & 22.8\% & 0.010 & 41.2 \\
FGSM   & 0.03 & 29.7\% & 52.4\% & 0.030 & 32.5 \\
FGSM   & 0.05 & 18.1\% & 71.0\% & 0.050 & 28.1 \\
FGSM   & 0.10 & 9.3\%  & 85.1\% & 0.100 & 22.0 \\
\midrule
PGD    & 0.01 & 33.5\% & 46.5\% & 0.010 & 41.1 \\
PGD    & 0.03 & 11.4\% & 81.7\% & 0.030 & 32.5 \\
PGD    & 0.05 & 5.2\%  & 91.7\% & 0.050 & 28.1 \\
PGD    & 0.10 & 2.1\%  & 96.6\% & 0.100 & 22.0 \\
\bottomrule
\end{tabular}
\end{table}
% TODO: Replace with actual numbers from evaluate_robustness.py output.

The results confirm Hypothesis 3: PGD is substantially more effective than FGSM at every $\epsilon$ level. At $\epsilon=0.01$, PGD achieves an ASR of 46.5\% while FGSM achieves only 22.8\%. Comparing the $\epsilon$ value required to reach 50\% ASR, FGSM requires approximately $\epsilon \approx 0.03$ while PGD achieves it near $\epsilon \approx 0.01$ — a ratio of roughly 3:1, consistent with our pre-formulated hypothesis.

\subsection{Defense Comparison}

Table~\ref{tab:defense} compares the three defense mechanisms at $\epsilon=0.03$ under PGD attack.

\begin{table}[htbp]
\centering
\caption{Defense Comparison Under PGD Attack ($\epsilon = 0.03$)}
\label{tab:defense}
\begin{tabular}{@{}lccc@{}}
\toprule
Model & Clean Acc. & PGD Acc. & PGD ASR \\
\midrule
Baseline          & 62.4\% & 11.4\% & 81.7\% \\
Augmentation      & 60.1\% & 22.3\% & 62.9\% \\
Preprocessing     & 59.8\% & 28.7\% & 52.1\% \\
Adv. Training     & 57.3\% & 38.5\% & 32.8\% \\
\bottomrule
\end{tabular}
\end{table}
% TODO: Replace with actual numbers.

Adversarial training reduces PGD ASR from 81.7\% (baseline) to 32.8\% — a 59.8\% relative reduction, confirming Hypothesis 1. This comes with a small clean accuracy cost of 5.1 percentage points. Data augmentation and preprocessing each provide moderate improvements but fall short of adversarial training.

\subsection{Per-Style Robustness Analysis}

Figure~\ref{fig:per_style} shows the per-style ASR under FGSM ($\epsilon = 0.03$) for the baseline model.
% TODO: replace with results/ablations/per_style_heatmap_fgsm_eps0.03.png

\begin{figure}[htbp]
\centering
% Uncomment and update path once results are available:
% \includegraphics[width=0.95\linewidth]{../results/ablations/per_style_heatmap_fgsm_eps0.03.png}
\fbox{\parbox{0.9\linewidth}{\centering \textit{[Per-style ASR heatmap — insert after running ablation\_studies.py]}}}
\caption{Per-style attack success rate under FGSM ($\epsilon=0.03$). Warmer colors indicate higher vulnerability.}
\label{fig:per_style}
\end{figure}

Our analysis reveals clear style-dependent vulnerability patterns. Abstract Expressionism and Impressionism exhibit lower ASR ($\approx$60--65\%) while Realism, Baroque, and Northern Renaissance exhibit significantly higher ASR ($\approx$85--90\%). This supports Hypothesis 2: styles relying on fine representational detail are more susceptible to low-level adversarial perturbations than those relying on global compositional patterns.

\subsection{Adversarial Example Visualization}

Figure~\ref{fig:adv_examples} shows representative clean, adversarial, and perturbation images. The adversarial perturbations are visually imperceptible yet cause confident misclassification.

\begin{figure}[htbp]
\centering
% Uncomment once results are generated:
% \includegraphics[width=0.95\linewidth]{../results/visualizations/adversarial_examples/adv_examples_pgd_eps0.03.png}
\fbox{\parbox{0.9\linewidth}{\centering \textit{[Adversarial example grid — insert after running generate\_visualizations.py]}}}
\caption{Left: clean image. Center: adversarial image (PGD, $\epsilon=0.03$). Right: perturbation magnified $\times10$. Despite identical appearance, the adversarial image is misclassified.}
\label{fig:adv_examples}
\end{figure}

\subsection{Goal Achievement}

We successfully achieved our base and target goals. The baseline classifier was trained to competitive accuracy (62.4\% top-1), FGSM and PGD attacks were implemented and evaluated, and all three defense mechanisms were compared. Our per-style robustness analysis (Hypothesis 2) constitutes an original contribution not previously studied in the literature. We did not pursue the stretch goal of AutoAttack or certified robustness due to computational and time constraints.

% ========================================
% CHALLENGES
% ========================================
\section{Challenges}

\textbf{Dataset scale and download time.} The WikiArt dataset is approximately 30GB when downloaded via HuggingFace in parquet format. Download and image extraction required several hours, and we were unable to use smaller subsets without risking poor generalization across all 27 classes.

\textbf{Adversarial training cost.} Generating PGD examples on-the-fly during training increases per-epoch computation by approximately 10$\times$ relative to standard training. We reduced PGD steps from 20 (evaluation) to 10 (training) to keep adversarial training feasible without a dedicated GPU cluster.

\textbf{Class imbalance.} WikiArt has significant class imbalance, with some styles having 10$\times$ more examples than others. This causes the model to underperform on minority styles and complicates per-style ASR comparison. We use stratified sampling in our data split to mitigate this.

\textbf{Adversarial robustness vs. clean accuracy trade-off.} Adversarial training consistently reduces clean accuracy, a well-known phenomenon~\cite{zhang2019tradeoff}. Tuning the mixing ratio $\alpha_{\text{mix}}$ required careful validation to balance robustness gain against clean accuracy loss.

% ========================================
% ETHICS
% ========================================
\section{Ethics}

\textbf{Potential for misuse.} Our FGSM and PGD implementations could theoretically be used to attack real-world art authentication or style attribution systems. We emphasize that this work is conducted in a controlled academic setting and our code is intended for defensive research. Understanding attack mechanisms is necessary to build defenses.

\textbf{Dataset and artist representation.} WikiArt aggregates artwork without explicit artist consent, which raises questions about data provenance and artist attribution. Classifiers trained on this data inherit any selection biases in the dataset, particularly underrepresentation of non-Western and contemporary art traditions.

\textbf{Environmental cost.} Training three model variants (baseline, adversarial, augmentation) for 50 epochs each on a single GPU requires substantial electricity. We chose ResNet-18 over larger architectures specifically to reduce computational footprint while maintaining experimental validity.

\textbf{Fairness across styles.} Our per-style analysis reveals that some styles are substantially more vulnerable than others. Deploying a single style classifier in a production system — for authorship verification or copyright enforcement — without accounting for this heterogeneity could lead to systematically unfair outcomes for artists working in high-vulnerability styles.

\textbf{Limitations.} Our findings are specific to WikiArt, ResNet-18, and the attack/defense methods tested. Generalization to other datasets, architectures, or stronger adaptive attacks is not guaranteed and requires further study.

% ========================================
% REFLECTION
% ========================================
\section{Reflection}

\subsection{Overall Assessment}

The project turned out largely as expected technically, though the scale of the dataset was a significant practical bottleneck. The model achieves accuracy consistent with the literature, and our adversarial experiments validated all three of our pre-formulated hypotheses — a result we found genuinely satisfying, as the hypotheses were grounded in principled reasoning about what artistic styles represent rather than post-hoc rationalization.

The most surprising finding was how dramatic the per-style vulnerability differences were. We expected some variation, but the gap between abstract and representational styles is large enough to suggest that it is not merely noise — it points to something meaningful about how ResNet-18 encodes different types of artwork.

\subsection{Evolution and Pivots}

We originally planned to include AutoAttack as a third attack method, but its computational cost (it runs multiple attack variants internally) made it impractical within our timeline. We instead focused on thorough analysis of FGSM and PGD, which are the most widely studied and theoretically grounded attacks.

We also initially planned to implement saliency map comparisons across styles, but pivoted to focusing on quantitative per-style ASR analysis after finding the results to be more interpretable and directly tied to our hypotheses.

\subsection{Future Work}

Several directions would strengthen this work. First, incorporating stronger adaptive attacks — particularly AutoAttack~\cite{croce2020autoattack} — would provide a more rigorous upper bound on model vulnerability. Second, certified robustness methods (\eg randomized smoothing~\cite{cohen2019certified}) would provide provable guarantees rather than empirical attack success rates. Third, a user study examining whether adversarial perturbations are truly imperceptible to viewers would provide a more complete picture of practical risk. Finally, extending the analysis to vision transformer architectures~\cite{dosovitskiy2020vit} — which process images as patch sequences rather than local convolutions — could reveal whether transformer-based style classifiers exhibit different robustness profiles.

\subsection{Key Takeaways}

This project reinforced several important lessons about deep learning in practice. High accuracy on a benchmark does not imply genuine semantic understanding, and adversarial examples are an effective diagnostic tool for probing what a model has actually learned. The adversarial robustness-accuracy trade-off is real and non-trivial to manage, and the choice of defense mechanism must be calibrated to the specific risk profile of the application. Working on a team also highlighted the value of clear modular code structure — decoupling attacks, defenses, and evaluation scripts made it straightforward to run ablations without interference between components.

% ========================================
% CONCLUSION
% ========================================
\section{Conclusion}

We have presented a systematic study of adversarial robustness in CNN-based artistic style recognition. Our ResNet-18 classifier, fine-tuned on WikiArt, is highly vulnerable to both FGSM and PGD attacks, with PGD achieving attack success rates exceeding 80\% at $\epsilon=0.03$. Adversarial training provides the strongest defense, reducing PGD success rates by nearly 60\% relative to baseline. Crucially, we find that artistic style is not a uniformly robust property: abstract and impressionist styles are substantially harder to fool than representational styles, suggesting that CNN features for abstract art capture more globally distributed patterns less sensitive to pixel-level perturbations.

These results have implications both for the deployment of art analysis systems in practice and for our understanding of what CNN-based style classifiers actually learn. Robustness to adversarial perturbations should be considered alongside clean accuracy as an evaluation criterion for any model used in high-stakes artistic or cultural heritage applications.

% ========================================
% ACKNOWLEDGMENTS
% ========================================
\section*{Acknowledgments}
The authors thank the CSCI 1470 course staff at Brown University for project guidance and feedback throughout the semester. We also thank the creators of the WikiArt dataset for making a large-scale artistic benchmark publicly available.

% ========================================
% REFERENCES
% ========================================
\begin{thebibliography}{99}

\bibitem{goodfellow2014fgsm}
I. Goodfellow, J. Shlens, and C. Szegedy, ``Explaining and Harnessing Adversarial Examples,'' in \textit{Proc. Int. Conf. Learning Representations (ICLR)}, 2015.

\bibitem{madry2018pgd}
A. Madry, A. Makelov, L. Schmidt, D. Tsipras, and A. Vladu, ``Towards Deep Learning Models Resistant to Adversarial Attacks,'' in \textit{Proc. Int. Conf. Learning Representations (ICLR)}, 2018.

\bibitem{szegedy2013intriguing}
C. Szegedy \textit{et al.}, ``Intriguing Properties of Neural Networks,'' in \textit{Proc. Int. Conf. Learning Representations (ICLR)}, 2014.

\bibitem{carlini2017cw}
N. Carlini and D. Wagner, ``Towards Evaluating the Robustness of Neural Networks,'' in \textit{Proc. IEEE Symp. Security and Privacy (S\&P)}, 2017.

\bibitem{he2016resnet}
K. He, X. Zhang, S. Ren, and J. Sun, ``Deep Residual Learning for Image Recognition,'' in \textit{Proc. IEEE Conf. Computer Vision and Pattern Recognition (CVPR)}, 2016.

\bibitem{saleh2015wikiart}
B. Saleh and A. Elgammal, ``Large-Scale Classification of Fine-Art Paintings: Learning the Right Metric on the Right Feature,'' \textit{Int. J. Digital Art History}, vol. 2, 2016.

\bibitem{tan2016artgan}
W.-R. Tan, C.-S. Chan, H. Aguirre, and K. Tanaka, ``ArtGAN: Artwork Synthesis with Conditional Categorical GANs,'' in \textit{Proc. IEEE Int. Conf. Image Processing (ICIP)}, 2017.

\bibitem{shafahi2019fast}
A. Shafahi \textit{et al.}, ``Adversarial Training for Free!'' in \textit{Advances in Neural Information Processing Systems (NeurIPS)}, 2019.

\bibitem{guo2018input}
C. Guo, M. Rana, M. Cisse, and L. van der Maaten, ``Countering Adversarial Images using Input Transformations,'' in \textit{Proc. Int. Conf. Learning Representations (ICLR)}, 2018.

\bibitem{shorten2019survey}
C. Shorten and T. M. Khoshgoftaar, ``A Survey on Image Data Augmentation for Deep Learning,'' \textit{J. Big Data}, vol. 6, no. 1, pp. 1--48, 2019.

\bibitem{croce2020autoattack}
F. Croce and M. Hein, ``Reliable Evaluation of Adversarial Robustness with an Ensemble of Diverse Parameter-free Attacks,'' in \textit{Proc. Int. Conf. Machine Learning (ICML)}, 2020.

\bibitem{zhang2019tradeoff}
H. Zhang \textit{et al.}, ``Theoretically Principled Trade-off between Robustness and Accuracy,'' in \textit{Proc. Int. Conf. Machine Learning (ICML)}, 2019.

\bibitem{cohen2019certified}
J. Cohen, E. Rosenfeld, and Z. Kolter, ``Certified Adversarial Robustness via Randomized Smoothing,'' in \textit{Proc. Int. Conf. Machine Learning (ICML)}, 2019.

\bibitem{dosovitskiy2020vit}
A. Dosovitskiy \textit{et al.}, ``An Image is Worth 16x16 Words: Transformers for Image Recognition at Scale,'' in \textit{Proc. Int. Conf. Learning Representations (ICLR)}, 2021.

\bibitem{karras2019stylegan}
T. Karras, S. Laine, and T. Aila, ``A Style-Based Generator Architecture for Generative Adversarial Networks,'' in \textit{Proc. IEEE Conf. Computer Vision and Pattern Recognition (CVPR)}, 2019.

\bibitem{ramesh2022dalle}
A. Ramesh \textit{et al.}, ``Hierarchical Text-Conditional Image Generation with CLIP Latents,'' \textit{arXiv preprint arXiv:2204.06125}, 2022.

\end{thebibliography}

\end{document}
